\section{Setup}
\label{sec:setup}

Before you can use this workflow, you need to connect the three tools to each other once. If GitHub or Claude Code are new to you, do not worry: the steps below explain each action as you go, and you will only need to do this once. The daily loop described in the next section is much shorter.

One-time setup takes about ten minutes. You need an Overleaf Premium account (GitHub sync is a paid feature), a free GitHub account, and a machine with Node.js and Git installed.

\subsection{Link Overleaf to GitHub}

In Overleaf, open your project, then:

\begin{enumerate}[leftmargin=*]
  \item \textbf{Menu} $\rightarrow$ \textbf{GitHub} $\rightarrow$ \textbf{Create a new repository on GitHub}, or link an existing one.
  \item Authorise Overleaf to access the repository.
  \item Overleaf performs an initial push. You now have a repository with your \verb|main.tex|, \verb|.bib| files, figures, etc.
\end{enumerate}

From this point on, Overleaf will not auto-sync. You push from Overleaf (\textbf{Menu $\rightarrow$ GitHub $\rightarrow$ Push}) and pull from Overleaf (\textbf{Menu $\rightarrow$ GitHub $\rightarrow$ Pull}) explicitly.

\subsection{Clone the repository locally}

\begin{lstlisting}[style=shell]
$ git clone git@github.com:<you>/<your-paper>.git
$ cd <your-paper>
\end{lstlisting}

\subsection{Install Claude Code}

Claude Code is distributed via npm:

\begin{lstlisting}[style=shell]
$ npm install -g @anthropic-ai/claude-code
$ claude --version
\end{lstlisting}

\begin{tipbox}{First time in a terminal?}
Open the Terminal application on your computer (search ``Terminal'' on Mac or Windows). You only need to run these two commands once. The Claude Code documentation at the URL below walks you through authentication step by step.
\end{tipbox}

\noindent See \url{https://docs.claude.com/en/docs/claude-code} for the current installation instructions and authentication steps.

\subsection{Start a session in the project}

\begin{lstlisting}[style=shell]
$ cd <your-paper>
$ claude
\end{lstlisting}

\noindent On first run, Claude Code will walk you through authentication. Subsequent sessions start immediately.

\subsection{Add a \texttt{CLAUDE.md}}

A \verb|CLAUDE.md| at the repository root is Claude's memory across sessions. It is read automatically at the start of every session and anchors the model to your conventions so you do not have to repeat them.

A minimal template for a paper:

\begin{lstlisting}[style=tex]
# CLAUDE.md

## What this repo is
A LaTeX paper: "<working title>". Target venue: <venue, page limit>.

## Style
- British English.
- Past tense for methods; present for claims.
- Oxford comma; en-dash for ranges; em-dash for asides (no spaces).
- Defined terms are italicised on first use, roman thereafter.

## Citations
- `\citet{key}` when the author is a noun ("Smith et al. (2023) showed...").
- `\citep{key}` when parenthetical.
- Every `\cite` key must exist in `references.bib`. Never invent keys.

## File layout
- `main.tex` - preamble and `\input{}` of each section.
- `sections/*.tex` - one file per section.
- `references.bib` - single bibliography.
- `figures/` - source figures (PDF or TikZ).

## Do not commit
Build artefacts: `*.aux *.log *.out *.toc *.bbl *.blg *.fls
*.fdb_latexmk *.synctex.gz`. Compiled `main.pdf` is also ignored.

## Commit style
Imperative mood, <= 72 chars, describe the "why".
Example: "Tighten intro to fit 400-word limit for reviewer R1".
\end{lstlisting}

This file is both a working configuration and a template you can copy into your own repositories.
